\ssr{Введение}

Цели данной лабораторной работы:
\begin{enumerate}
	\item изучение системных средств обмена данными между процессами в ОС GNU/LINUX;
	\item получение практических навыков использования механизмов межпроцессного взаимодействия.
\end{enumerate}

\ssr{Выполнение}

\ssection{Задание 1}

Задание по варианту:

Реализовать перемножение двух матриц с помощью нескольких процессов: каждый процесс выполняет перемножение строки первой матрицы на столбец второй (в соответствии с правилом умножения матриц). При необходимости процессам, выполняющим умножение, может быть отправлено несколько заданий.

Код программы представлен в листинге~\ref{task1}.

% gcc communications.c -o communications -lrt -pthread

\lstinputlisting[
	style=cstyle,
	label={task1},
	caption={Программа по варианту}
]{communications.c}

\begin{figure}[!htbp]
	\centering
	\includegraphics[width=0.6\linewidth]{"Screenshot From 2026-11-16 12-59-50.png"}
	\caption{Демонстрация работы}
\end{figure}
\FloatBarrier

\begin{figure}[!htbp]
	\centering
	\includegraphics[width=0.6\linewidth]{"os_lab9.png"}
	\caption{Схема программы}
\end{figure}
\FloatBarrier

\ssection{Дополнительное задание}

Анализ выполнения задачи по варианту для различных средств межпроцессного взаимодействия представлен в таблице~\ref{comp}.

\begin{table}[!htbp]
	\caption{Анализ выполнения задачи}
		\label{comp}
	\begin{tabularx}{\textwidth}{|X|X|X|X|}
		\hline
		\textbf{Название} & \textbf{Время выполнения} & \textbf{Объем дополнительной используемой памяти} & \textbf{Количество сообщений между процессами} \\\hline
		Каналы & 0.000867с & 256КБ & 14 \\\hline
		% 2 * NUM_WORKERS * (4 * 16)
		Разделяемая память & 0.000869с & 4КБ & 8-10 \\\hline
		% sizeof(SharedData)
		Сокеты & 0.001315с & 624КБ & 24-27 \\\hline
		% (1 + NUM_WORKERS) * 208
	\end{tabularx}
\end{table}

Разделяемая память -- самый быстрый механизм для данной задачи: нет копирования данных через ядро, обмен идет через общую область. Однако требует явной синхронизации (атомарные операции или семафоры), что усложняет код и создает риск гонок.

Каналы -- промежуточный вариант: данные копируются через буферы ядра, но синхронизация встроена (блокирующие read/write). Для большого числа ведомых процессов требуется select/poll для мультиплексирования, что усложняет программу.

Сокеты -- самый медленный на одной машине из-за накладных расходов на формирование и разбор адресов, проверку прав, буферизации, но самый  универсальный: тот же код можно использовать для сетевого взаимодействия (заменив AF\_UNIX на AF\_INET), и они дают естественную адресацию клиентов.

По объему памяти разделяемая память выигрывает: один сегмент на всех, тогда как каналы и сокеты требуют отдельных буферов ядра на каждое соединение.
По количеству сообщений каналы и разделяемая память экономнее сокетов, так как в сокетах добавляются служебные запросы (получение матриц, подтверждения).

Таким образом, для задачи перемножения матриц на одной машине оптимальна разделяемая память (с семафорами для синхронизации); если нужна простота и надежность -- каналы; если в перспективе -- сокеты.

\ssr{Вывод}

В ходе лабораторной работы были:  
\begin{itemize}
	\item изучены системные средства обмена данными между процессами в ОС GNU/LINUX;
	\item получены практическиу навыки использования механизмов межпроцессного взаимодействия.
\end{itemize}

\ssr{Контрольные вопросы}

\begin{enumerate}
	\item \textit{Дайте определение механизмам межпроцессного взаимодействия (IPC)? Поясните, почему возникает задача коммуникации между различными процессами.}\newline 
	Это механизмы, позволяющие процессам обмениваться данными и синхронизировать действия (каналы, разделяемая память, сокеты, очереди сообщений, семафоры). 
	Задача коммуникации возникает, когда один процесс должен передать данные другому (например, ведущий-ведомый, клиент-сервер), а адресные пространства процессов изолированы.
	\item \textit{В чём различие между использованием неименованных и именованных каналов?}\newline 
	Неименованные каналы доступны только создавшему процессу и его потомкам (через наследование дескрипторов), не имеют имени в ФС. 
	Именованные (FIFO) существуют как файл в ФС, к ним могут подключаться произвольные процессы, знающие имя и имеющие права.
	\item \textit{Почему при использовании разделяемой памяти необходимы дополнительные механизмы синхронизации процессов?}\newline
	Разделяемая память не содержит встроенной синхронизации: несколько процессов могут одновременно читать/писать одну область, что приводит к гонкам.
	\item \textit{В чем разница между флагами MAP\_SHARED и MAP\_PRIVATE при вызове mmap()?}\newline 
	MAP\_SHARED -- изменения видны всем процессам, отобразившим тот же объект, и записываются в файл. 
	MAP\_PRIVATE -- флаг copy-on-write: изменения видны только данному процессу и не влияют на файл.
	\item \textit{В чем разница между сокетами в домене AF\_UNIX и домене AF\_INET?}\newline 
	AF\_UNIX -- сокеты в файловом пространстве имен (локальные, через файл-сокет), быстрее, без сетевого стека. 
	AF\_INET -- сетевые сокеты (IP-адрес + порт), работают по сети, но имеют накладные расходы.
	\item \textit{Какую роль выполняет функция bind() в программе-сервере?}\newline 
	Она связывает сокет с локальным адресом (именем файла для AF\_UNIX или IP-портом для AF\_INET), после чего сервер становится доступен клиентам по этому адресу.
	\item \textit{Какие функции в серверном, а какие в клиентском процессе необходимо выполнить для организации коммуникации, если сокет создан функцией: int sock = socket(AF\_UNIX, SOCK\_STREAM, 0); ?}\newline 
	Сервер -- socket(), bind(), listen(), accept(), read()/write(), close(), unlink(). 
	Клиент -- socket(), connect(), read()/write(), close().
	\item \textit{Какие преимущества и недостатки у использования сокетов по сравнению с каналами для межпроцессного взаимодействия на одной машине?}\newline 
	Преимущества: универсальность (локально и по сети), адресация клиентов, стандартный API. 
	Недостатки: больше накладных расходов, сложнее код, медленнее каналов на одной машине.
	\item \textit{Какой механизм IPC наиболее подходит для передачи больших объемов данных и почему?}\newline 
	Для больших объемов данных лучше разделяемая память, т.к. нет копирования через ядро, минимальные накладные расходы; при необходимости -- mmap файла.
	\item \textit{Что такое «файловые дескрипторы»? Какие файловые дескрипторы и какими функциями создаются для рассмотренных методов IPC?}\newline 
	Файловый дескриптор -- целое число, индекс в таблице открытых файлов процесса. 
	Для IPC: pipe() -> 2 дескриптора; mkfifo() + open() -> дескриптор FIFO; shm\_open() -> дескриптор объекта shm; socket() -> дескриптор сокета; accept() -> дескриптор соединения.
	\item \textit{Как работает системный вызов poll() и чем он отличается от select()}\newline 
	poll() -- ожидание готовности набора дескрипторов, использует массив структур pollfd (без ограничения FD\_SETSIZE, в отличие от select), не изменяет наборы, поддерживает больше типов событий (POLLIN, POLLOUT, POLLERR, POLLHUP и др.). 
	select() использует битовые маски и ограничен FD\_SETSIZE (обычно 1024).
	poll() масштабируется лучше при большом числе дескрипторов.
\end{enumerate}
